\documentclass[12pt]{article}
\usepackage{styles/cwilliams-standard}
\usepackage{tablefootnote}
\usepackage{multirow}
\usepackage{listings}
\lstset{
    basicstyle=\ttfamily\small,
    numbers=left,
    numberstyle=\tiny\color{gray},
    numbersep=8pt,
    breaklines=true,
    frame=single,
    rulecolor=\color{gray!40},
    backgroundcolor=\color{gray!6},
    keywordstyle=\color{blue!70},
    stringstyle=\color{orange!80!black},
    showstringspaces=false,
    language=Python,
}

\setclass{\NLP}
\settitle{Homework 2: Affect, Perspective, and Subective Judgement}

\definecolor{tableheader}{rgb}{0.18, 0.32, 0.55}
\definecolor{rowalt}{rgb}{0.93, 0.95, 0.98}
\definecolor{modelheader}{rgb}{0.05, 0.42, 0.42}
\definecolor{modelhighlight}{rgb}{1.0, 0.88, 0.50}

\begin{document}

\maketitlepage

\tableofcontents

\begin{important}{Table Colors}
    Each task will have different colored tables, but the functions of highlights will remain the 
    same the entire time. I did this to make quickly scrolling through the document easier to 
    quickly know where you are in the report. Light blue is Task 1, teal and amber is Task 2, etc..
\end{important}

\section{Affective Analysis}

\begin{enumerate}
    \item \textbf{Leveraging lexicons provided by Mohammed et al.}
    
    To ``leverage'' the lexicons provided, I imported the NRC-VAD-Lexicon-v2.1 from 
    this path: \texttt{NRC-VAD-Lexicon-v2.1/Unigrams/unigrams-NRC-VAD-Lexicon-v2.1.txt}. 
    
    I loaded this into a dictionary where the term is the key, and the valence, arousal, dominance (VAD) are values. 
    This allows for a faster lookup (O(1)) and quicker processing down the line. Given that the 
    downloaded files already have VAD provided, there was no computation that 
    really needed to be done.

    \item \textbf{Relevant Sentence Extraction}
    
    Given the three focus words (spy, stalk, stealth), I went through all sentences in the dataset
    \texttt{reviewer.csv} and used the spaCy \texttt{en\_core\_web\_sm} package to seperate each word in each of the sentences. 
    Looping through the split sentences, I then check if the sentence contains one of the three focus words. 
    If it does contain the focus word, save the entire sentence along with the reviewer type, add that to a list 
    as a tuple, and the move on. 

    \begin{table}[H]
        \centering
        \renewcommand{\arraystretch}{1.4}
        \begin{tabular}{p{10cm} p{3cm}}
            \toprule
            \rowcolor{tableheader}
            \color{white}\textbf{Sentence} & \color{white}\textbf{Reviewer} \\
            \midrule
            \rowcolor{rowalt} I love that you can take a picture and am constantly spying on coworkers at my desk from my laptop. & abuser \\
            This is also good to spy on someone. & third\_party \\
            \rowcolor{rowalt} i can spy on my child whenever i want its amazing he cant go anywhere without me knowing look & abuser \\
            Great app for families keeping track of their kids, friends or if they just wanna be creepy stalkers. & abuser \\
            \rowcolor{rowalt} Another way of spying on people. & abuser \\
            \bottomrule
        \end{tabular}
        \caption{Example Sentences Containing Focus Words}
        \label{tab:sentences}
    \end{table}

    \item \textbf{Adjective Extraction}
    
    For this subtask, I went through all sentences in the \texttt{reviewer.csv} and used the spaCy 
    \texttt{en\_core\_web\_sm} package again to POS tag every word. Once every word is tagged, go through each 
    sentence and then each word, and if the word's POS tag is ``\texttt{ADJ}'', then add that word and it's 
    corresponding VAD scores to a dictionary where the key is the reviewer type, and the 
    values are a list of dictionary's that have the ADJ's as keys and the values are their VAD scores.

\begin{lstlisting}
{   'third_party': [
        {'word': 'good', 
        'valence': 0.876, 'arousal': -0.264, 'dominance': 0.068},
        ...
    ],
    'abuser': [
        {'word': 'amazing',
         'valence': 0.812, 'arousal': 0.674, 'dominance': 0.698}, 
        ...
    ],
    'victim': [
        {'word': 'teen',
         'valence': 0.458, 'arousal': 0.062, 'dominance': 0.136}, 
         ...
    ]
}
\end{lstlisting}

    \item \textbf{Average scores of adjectives for all dimensions}
    
    Using the dictionary created in the previous subtask, I go by key (reviewer type) and get all scores and average them
    across the VAD domains. This is rather straight forward, and here are the scores that I got for this task.

    \begin{table}[H]
        \centering
        \renewcommand{\arraystretch}{1.4}
        \begin{tabular}{p{4cm} p{3cm} p{3cm} p{3cm}}
            \toprule
            \rowcolor{tableheader}
            \color{white}\textbf{Reviewer} & \color{white}\textbf{Valence} & \color{white}\textbf{Arousal} & \color{white}\textbf{Dominance} \\
            \midrule
            \rowcolor{rowalt} third\_party & 0.186 & -0.006 & 0.055 \\
                              abuser & 0.498 & -0.005 & 0.237 \\
            \rowcolor{rowalt} victim & 0.100 & 0.055 & 0.034 \\
            \bottomrule
        \end{tabular}
        \caption{Average VAD Scores by Reviewer Type}
        \label{tab:vad-averages}
    \end{table}

    \item \textbf{Compare dimensions}
    
    As per the table in the previous subtask, we can see there are values per reviewer type that 
    describe emotions and intensity of emotions. 

    \begin{table}[H]
        \centering
        \renewcommand{\arraystretch}{1.4}
        \begin{tabular}{p{4cm} p{3cm} p{3cm} p{3cm}}
            \toprule
            \rowcolor{tableheader}
            \color{white}\textbf{Reviewer} & \color{white}\textbf{Valence} & \color{white}\textbf{Arousal} & \color{white}\textbf{Dominance} \\
            \midrule
            \rowcolor{rowalt} third\_party & 0.186 & \cellcolor{red!13}\textbf{-0.006} & 0.055 \\
                              abuser & \cellcolor{green!13}\textbf{0.498} & -0.005 & \cellcolor{green!13}\textbf{0.237} \\
            \rowcolor{rowalt} victim & \cellcolor{red!13}\textbf{0.100} & \cellcolor{green!13}\textbf{0.055} & \cellcolor{red!13}\textbf{0.034} \\
            \bottomrule
        \end{tabular}
        \caption{Average VAD Scores by Reviewer Type w/ Highest Value Highlighted. 
        The scores highlighted red are the lowest for that VAD domain. 
        The scores highlighted green are the highest values for that VAD domain}
        \label{tab:vad-averages-highlighted}
    \end{table}

    Abusers appear to have very positive feelings towards ($0.498$ valence) the apps they review and feel a higher 
    degree of control with a dominance of $0.237$ on average when reviewing about the app, showing that they feel they 
    have control over who they are ``stalking'' or abusing. Conversely, victims have less positive feelings (although not negative) 
    about the apps but feel more intensely about the apps with an arousal of $0.055$ on average and feel much less in control than 
    the abusers with a dominance of $0.034$. These two go hand in hand and it makes sense overall that the third\_party reviwers do not have particular strong 
    feelings towards this app as they are not the ones using the apps. 

    Third\_party reviewers have more positive feelings about this than victims which does make sense given that 
    they are not experiencing the effects of these apps and are only commenting on them. And, they seem to feel 
    more in control than victims which also makes sense.

    Ultimately, it appears that abusers feel the most in control in relation to these apps, victims feel the 
    least positive and in control, and third\_party is in between.

\end{enumerate}

\section{Joint Modeling for Severity and Convincingness}

For this task, I selected BERT and DistilBERT for the two models that I will be comparing. 

\begin{table}[H]
    \centering
    \renewcommand{\arraystretch}{1.4}
    \begin{tabular}{l l c c c c}
        \toprule
        \rowcolor{modelheader}
        \color{white}\textbf{Model} & \color{white}\textbf{Type} & \color{white}\textbf{Precision} & \color{white}\textbf{Recall} & \color{white}\textbf{F1} & \color{white}\textbf{Macro F1} \\
        \midrule
        \multirow{2}{*}{BERT}       & Severity       & \cellcolor{modelhighlight}\textbf{0.7480} & \cellcolor{modelhighlight}\textbf{0.7401} & \cellcolor{modelhighlight}\textbf{0.7336} & \multirow{2}{*}{0.6755} \\
                                     & Convincingness & 0.6392 & 0.6378 & 0.6174 & \\
        \cmidrule{1-6}
                          \multirow{2}{*}{DistilBERT} & Severity       & 0.7140 & 0.7280 & 0.7190 & \multirow{2}{*}{0.6817} \\
                                                       & Convincingness & \cellcolor{modelhighlight}\textbf{0.6978} & \cellcolor{modelhighlight}\textbf{0.6543} & \cellcolor{modelhighlight}\textbf{0.6444} & \\
        \bottomrule
    \end{tabular}
    \caption{BERT vs.\ DistilBERT Classification Performance. Amber highlighted cells show the best
    performance for that type across models. One cell could be read as ``BERT has the best score for
    precision in severity across both models''}
    \label{tab:model-comparison}
\end{table}

\section{Role-Aware LLM Analysis}

\subsection{Zero Shot \& Few-Shot Detection}

\subsection{Comparing Performance}

\section{Code}

% #TODO
\begin{aside}{Provided Code}
    The code is provided with this submission on Canvas, but in case something is not correct 
    or does not work, \href{github}{here} is the link to the repo on Github.
\end{aside}

\end{document}